\documentclass[12pt,a4paper]{article}
\usepackage[margin=1in]{geometry}
\usepackage{amsmath,amssymb}
\usepackage{booktabs}
\usepackage{graphicx}
\usepackage{hyperref}
\usepackage{setspace}
\usepackage{caption}
\usepackage{array}
\usepackage{float}
\usepackage{xcolor}
\usepackage{listings}

\setstretch{1.15}

\definecolor{codegreen}{rgb}{0,0.6,0}
\definecolor{codegray}{rgb}{0.5,0.5,0.5}
\definecolor{codepurple}{rgb}{0.58,0,0.82}
\definecolor{backcolour}{rgb}{0.95,0.95,0.92}

\lstdefinestyle{mystyle}{
    backgroundcolor=\color{backcolour},
    commentstyle=\color{codegreen},
    keywordstyle=\color{magenta},
    numberstyle=\tiny\color{codegray},
    stringstyle=\color{codepurple},
    basicstyle=\ttfamily\footnotesize,
    breakatwhitespace=false,
    breaklines=true,
    captionpos=b,
    keepspaces=true,
    numbers=left,
    numbersep=5pt,
    showspaces=false,
    showstringspaces=false,
    showtabs=false,
    tabsize=2
}

\lstset{style=mystyle}

\title{\textbf{ITI123 Final Project Report}\\\vspace{0.3cm}AI-Based Badminton Shot Classification:\\From Pose-Only Failure to Video-Based Success}
\author{\textbf{Paul George Karippaparambil}\\8031408E}
\date{ITI123\\February 06, 2026}

\begin{document}
\maketitle
\thispagestyle{empty}

\begin{abstract}
This project explores the application of deep learning for automated badminton shot classification. We implemented and evaluated three distinct approaches: (1) pose-only skeleton-based models achieving 50\% accuracy on Clear vs Smash binary classification, (2) video-based ResNet18+BiLSTM achieving 74.6\% accuracy on 5-class shot type classification, and (3) an end-user application (Shot Coach) providing real-time shot analysis with video metadata tracking. Key findings demonstrate that visual appearance features significantly outperform skeletal pose features for badminton shot recognition. The final system processes video clips in 5-10 seconds and provides actionable feedback through an intuitive web interface. This report documents the complete development journey including failed approaches, lessons learned, evaluation results, and practical improvements.
\end{abstract}

\newpage
\tableofcontents
\newpage

\section{Introduction}

\subsection{Motivation}

As a recreational badminton player, obtaining technique feedback without a professional coach presents significant challenges. Traditional coaching requires scheduling, travel, and financial investment. This project investigates whether artificial intelligence can provide automated shot classification and technique analysis, making badminton coaching more accessible.

\subsection{Research Questions}

This project addresses three core questions:

\begin{enumerate}
\item \textbf{Phase 1 (Pose-Only):} Can skeletal keypoint data alone distinguish between Clear and Smash shots?
\item \textbf{Phase 2 (Video-Based):} Can visual appearance features improve classification across all 5 shot types?
\item \textbf{Phase 3 (Application):} Can the trained model be deployed in a practical user-facing application?
\end{enumerate}

\subsection{Shot Types}

The project classifies five badminton shot types:

\begin{table}[H]
\centering
\caption{Badminton Shot Types}
\begin{tabular}{lp{8cm}}
\toprule
Shot Type & Description \\
\midrule
\textbf{Clear} & High defensive shot to the back of the court, used to push opponent back and buy time \\
\textbf{Drive} & Fast, flat shot parallel to the ground, used for quick attacks \\
\textbf{Drop} & Soft shot that barely clears the net, forcing opponent to move forward \\
\textbf{Lift} & Defensive shot hit high to the back, similar to Clear but from defensive position \\
\textbf{Smash} & Powerful attacking shot hit downward, primary attacking shot in badminton \\
\bottomrule
\end{tabular}
\end{table}

\subsection{Dataset}

All experiments use the ShuttleSet dataset \cite{ShuttleSet}, a comprehensive collection of annotated badminton match videos:

\begin{itemize}
\item Professional match footage from international tournaments
\item Frame-level annotations with shot types, player positions, and timing
\item 27,374 total annotated shots across 19 shot type categories
\item 22,302 samples after filtering to 5 primary shot types
\item 40 different matches providing diverse playing styles
\end{itemize}

\section{Phase 1: Pose-Only Classification (Failed Approach)}

\subsection{Methodology}

The initial approach tested whether body pose alone could distinguish between Clear and Smash shots---two overhead shots with similar preparation but different execution.

\subsubsection{Data Preprocessing}

\textbf{Shot Selection:}
\begin{itemize}
\item Started with 4,983 Clear and Smash clips
\item Removed 301 backhand clips (different biomechanics)
\item Final dataset: 4,682 forehand clips
\item Class distribution: 2,371 Clear (50.6\%), 2,311 Smash (49.4\%)
\end{itemize}

\textbf{Pose Extraction:}

Used Google MediaPipe Pose \cite{MediaPipe} with highest accuracy settings:
\begin{itemize}
\item Model complexity: 2 (most accurate)
\item 33 body keypoints per frame (x, y, z coordinates)
\item Key landmarks: shoulders, elbows, wrists, hips, knees, ankles, nose
\item Success rate: 99.4\% (4,655 clips with valid pose data)
\end{itemize}

\subsubsection{Feature Engineering}

To capture shot-specific biomechanics, I engineered 60 features per frame across multiple categories:

\begin{table}[H]
\centering
\caption{Pose Feature Categories}
\begin{tabular}{llc}
\toprule
Category & Features & Count \\
\midrule
Arm Extension & Shoulder-to-wrist distance & 2 \\
Depth Features & Wrist/elbow depth relative to shoulder/hip & 4 \\
Height Features & Wrist/elbow height relative to shoulder/head & 4 \\
Joint Angles & Elbow angle, shoulder angle & 2 \\
Body Posture & Torso lean, shoulder rotation & 4 \\
Wrist/Forearm & Vertical angle, pitch, pronation, reach & 9 \\
Velocities & First derivative of all spatial features & 30 \\
Other & Hip width, stance, balance metrics & 5 \\
\midrule
\textbf{Total} & & \textbf{60} \\
\bottomrule
\end{tabular}
\end{table}

\textbf{Key Hypothesis:} Wrist angle would discriminate shots. Expected Clear to use extended wrist (120--150$^\circ$) and Smash to use flexed wrist (30--60$^\circ$).

\subsubsection{Data Splitting}

Implemented group-stratified splitting to prevent data leakage:

\begin{itemize}
\item All clips from the same match stay together (avoid learning match-specific patterns)
\item Stratified by shot type to maintain class balance
\item Train: 3,554 clips (27 matches), Val: 426 clips (5 matches), Test: 675 clips (8 matches)
\end{itemize}

\begin{table}[H]
\centering
\caption{Pose-Only Dataset Split}
\begin{tabular}{lcccc}
\toprule
Split & Clips & Matches & \% Clear & \% Smash \\
\midrule
Train & 3,554 & 27 & 50.45 & 49.55 \\
Validation & 426 & 5 & 50.70 & 49.30 \\
Test & 675 & 8 & 54.22 & 45.78 \\
\midrule
Total & 4,655 & 40 & 50.62 & 49.38 \\
\bottomrule
\end{tabular}
\end{table}

\subsubsection{Model Architectures}

Tested four baseline architectures for temporal sequence modeling:

\begin{enumerate}
\item \textbf{ResNet1D}: 1D CNN with residual connections for pattern detection across motion sequences
\item \textbf{LSTM}: Long Short-Term Memory network capturing temporal dependencies
\item \textbf{BiLSTM}: Bidirectional LSTM processing sequences forward and backward
\item \textbf{GRU}: Gated Recurrent Unit, simpler and faster than LSTM
\end{enumerate}

All models used:
\begin{itemize}
\item Input shape: (50 timesteps, 60 features)
\item Padding masking to handle variable-length sequences
\item Adam optimizer (learning rate 0.001)
\item Binary crossentropy loss
\item Early stopping (patience 10 epochs)
\end{itemize}

\subsection{Results}

\begin{table}[H]
\centering
\caption{Pose-Only Baseline Results (Test Set)}
\begin{tabular}{lcccc}
\toprule
Model & Accuracy & Precision & Recall & F1 \\
\midrule
ResNet1D & 0.491 & 0.487 & 0.507 & 0.497 \\
LSTM & 0.457 & 0.459 & 0.446 & 0.452 \\
BiLSTM & 0.433 & 0.428 & 0.476 & 0.451 \\
GRU & 0.450 & 0.463 & 0.504 & 0.483 \\
\midrule
Random Baseline & 0.500 & 0.500 & 0.500 & 0.500 \\
\bottomrule
\end{tabular}
\end{table}

\textbf{Key Finding:} All models performed at or below random guessing (50\% accuracy).

\subsubsection{Enhanced Models}

Tested sophisticated architectures to verify that the problem was data quality, not model capacity:

\begin{itemize}
\item Attention mechanisms (Bahdanau attention, self-attention)
\item Transformer architecture with positional encoding
\item Data augmentation (time shifts, scaling, noise injection, mixup)
\item Squeeze-and-Excitation blocks for channel recalibration
\item Label smoothing to prevent overconfidence
\end{itemize}

\begin{table}[H]
\centering
\caption{Enhanced Model Results (Test Set)}
\begin{tabular}{lcccc}
\toprule
Model & Accuracy & Precision & Recall & F1 \\
\midrule
ResNet1D + SE Block & 0.51 & 0.51 & 0.50 & 0.50 \\
LSTM + Attention & 0.50 & 0.50 & 0.51 & 0.50 \\
BiLSTM + Attention & 0.49 & 0.49 & 0.50 & 0.49 \\
Transformer & 0.50 & 0.50 & 0.49 & 0.50 \\
\midrule
Random Baseline & 0.50 & 0.50 & 0.50 & 0.50 \\
\bottomrule
\end{tabular}
\end{table}

\textbf{Conclusion:} No improvement with increased model complexity. Problem is inherent to the data representation, not model architecture.

\subsection{Root Cause Analysis}

\subsubsection{Identical Body Poses at Contact}

Analysis revealed that Clear and Smash have indistinguishable body poses at the moment of shuttlecock contact:

\begin{table}[H]
\centering
\caption{Body Measurements at Contact (Mean $\pm$ Std)}
\begin{tabular}{lcc}
\toprule
Feature & Clear & Smash \\
\midrule
Contact Height & 2.1m $\pm$ 0.3m & 2.0m $\pm$ 0.3m \\
Arm Extension & 94.2\% $\pm$ 8.1\% & 93.8\% $\pm$ 8.3\% \\
Shoulder Angle & 162.1$^\circ$ $\pm$ 12.4$^\circ$ & 161.8$^\circ$ $\pm$ 13.1$^\circ$ \\
Elbow Angle & 172.3$^\circ$ $\pm$ 9.2$^\circ$ & 171.9$^\circ$ $\pm$ 9.8$^\circ$ \\
Forearm Angle & 85.3$^\circ$ $\pm$ 45.2$^\circ$ & 84.2$^\circ$ $\pm$ 44.8$^\circ$ \\
\midrule
\textbf{Max Difference} & \multicolumn{2}{c}{1.1$^\circ$ (negligible)} \\
\bottomrule
\end{tabular}
\end{table}

The wrist angle hypothesis was disproven: expected 60--90$^\circ$ difference, observed 1.1$^\circ$ difference.

\subsubsection{Missing Critical Information}

Key discriminative features are absent from pose data:

\begin{enumerate}
\item \textbf{Racket face angle:} Clear angles racket upward (45--60$^\circ$), Smash angles downward (-30--0$^\circ$). MediaPipe doesn't track equipment. \textit{Most important feature---completely missing.}

\item \textbf{Shuttlecock trajectory:} Clears travel high and deep, Smashes travel steep and fast. Video clips end at contact. No trajectory information.

\item \textbf{Follow-through motion:} Clears continue upward swing, Smashes snap downward. Clips truncated at contact. No post-contact data.

\item \textbf{Racket head velocity:} Smashes have higher racket speeds ($>$300 km/h vs $<$200 km/h). Body velocity doesn't capture racket dynamics.
\end{enumerate}

\subsection{Lessons Learned}

\begin{enumerate}
\item \textbf{Domain knowledge matters:} Understanding sport biomechanics early would have revealed pose-only limitation
\item \textbf{Feature analysis before modeling:} Should have analyzed feature distributions before training
\item \textbf{Equipment tracking essential:} Sport analysis requires tracking sports equipment, not just body
\item \textbf{Temporal context needed:} Single-moment analysis insufficient for dynamic actions
\end{enumerate}

\section{Phase 2: Video-Based Classification (Successful Approach)}

\subsection{Motivation for Pivot}

Phase 1 failure motivated transition to video-based approach:
\begin{itemize}
\item Visual appearance captures racket orientation, ball trajectory
\item Temporal CNNs can learn discriminative patterns from pixels
\item Transfer learning from ImageNet provides strong feature extraction
\item Proven success in action recognition literature
\end{itemize}

\subsection{Architecture Design}

\subsubsection{Model: ResNet18 + Bidirectional LSTM}

\begin{figure}[H]
\centering
\begin{verbatim}
Video (2-3 seconds, 30 FPS)
    ↓
Sample 16 frames uniformly
    ↓
ResNet18 CNN (per frame)
    → Extract 512-dim spatial features
    ↓
Bidirectional LSTM (2 layers, 256 hidden units)
    → Capture temporal dependencies
    ↓
Fully Connected Layer
    → 5-class softmax output
\end{verbatim}
\caption{Video-Based Model Architecture}
\end{figure}

\textbf{Design Rationale:}
\begin{itemize}
\item \textbf{ResNet18 backbone:} Pretrained on ImageNet, provides strong spatial feature extraction with residual connections preventing vanishing gradients
\item \textbf{Bidirectional LSTM:} Processes sequence both forward and backward, capturing temporal patterns in both directions
\item \textbf{16 frames:} Balances temporal coverage (0.5s at 30 FPS) with computational efficiency
\item \textbf{Transfer learning:} Fine-tune on badminton domain while leveraging general visual knowledge
\end{itemize}

\subsubsection{Dataset Expansion}

Extended from 2-class to 5-class classification:

\begin{table}[H]
\centering
\caption{Video-Based Dataset Statistics}
\begin{tabular}{lrrc}
\toprule
Shot Type & Clips & Percentage & Difficulty \\
\midrule
Drop & 5,773 & 31.8\% & Easy \\
Lift & 5,230 & 28.8\% & Moderate \\
Smash & 3,872 & 21.3\% & Moderate \\
Clear & 2,662 & 14.6\% & Hard \\
Drive & 630 & 3.5\% & Very Hard \\
\midrule
\textbf{Total} & \textbf{22,302} & \textbf{100\%} & \\
\bottomrule
\end{tabular}
\end{table}

\textbf{Data Splits:}
\begin{itemize}
\item Train: 15,611 clips (70\%)
\item Validation: 2,208 clips (10\%)
\item Test: 4,483 clips (20\%)
\end{itemize}

\subsection{Training Configuration}

\subsubsection{Handling Class Imbalance}

Implemented Focal Loss \cite{FocalLoss} to address severe class imbalance:

\begin{equation}
FL(p_t) = -(1-p_t)^\gamma \log(p_t)
\end{equation}

where $\gamma = 2$ focuses training on hard examples by down-weighting easy classifications.

Class weights derived from inverse frequency:
\begin{equation}
w_i = \frac{N_{total}}{N_i \times C}
\end{equation}

\begin{table}[H]
\centering
\caption{Focal Loss Class Weights}
\begin{tabular}{lrrr}
\toprule
Shot Type & Samples & Weight & Relative Emphasis \\
\midrule
Drive & 630 & 5.67 & 5.67$\times$ \\
Clear & 2,662 & 1.68 & 1.68$\times$ \\
Smash & 3,872 & 1.21 & 1.21$\times$ \\
Lift & 5,230 & 0.92 & 0.92$\times$ \\
Drop & 5,773 & 0.84 & 0.84$\times$ \\
\bottomrule
\end{tabular}
\end{table}

\subsubsection{Hyperparameters}

\begin{table}[H]
\centering
\caption{Training Configuration}
\begin{tabular}{ll}
\toprule
Parameter & Value \\
\midrule
Optimizer & Adam \\
Learning Rate & 0.001 (initial) $\to$ 0.0001 (fine-tuning) \\
Batch Size & 64 \\
Loss Function & Focal Loss ($\gamma=2$, class weights) \\
Epochs & 50 (early stopping patience 15) \\
Frame Size & 224 $\times$ 224 pixels \\
Frames per Video & 16 (uniformly sampled) \\
Data Augmentation & Random horizontal flip, color jitter \\
Gradient Clipping & Max norm 1.0 \\
\bottomrule
\end{tabular}
\end{table}

\subsection{Results and Evaluation}

\subsubsection{Overall Performance}

\begin{table}[H]
\centering
\caption{Final Model Performance (Test Set)}
\begin{tabular}{lc}
\toprule
Metric & Value \\
\midrule
Test Accuracy & \textbf{74.62\%} \\
Macro F1-Score & 73.61\% \\
Weighted F1-Score & 74.87\% \\
Total Test Samples & 4,483 \\
Correct Predictions & 3,345 \\
Training Epochs & 44 (early stopped) \\
Best Validation Accuracy & 75.41\% \\
\bottomrule
\end{tabular}
\end{table}

\textbf{Comparison to Pose-Only:}
\begin{itemize}
\item Pose-only: 50\% accuracy (random guessing)
\item Video-based: 74.6\% accuracy
\item \textbf{Improvement: +24.6 percentage points (+49\% relative)}
\end{itemize}

\subsubsection{Per-Class Performance}

\begin{table}[H]
\centering
\caption{Classification Report by Shot Type}
\begin{tabular}{lcccc}
\toprule
Shot Type & Precision & Recall & F1-Score & Support \\
\midrule
Clear & 0.676 & \textbf{0.892} & 0.769 & 535 \\
Drive & 0.583 & 0.616 & 0.599 & 740 \\
Drop & \textbf{0.907} & 0.746 & 0.819 & 1,518 \\
Lift & 0.710 & 0.757 & 0.733 & 966 \\
Smash & 0.765 & 0.758 & 0.761 & 724 \\
\midrule
Macro Avg & 0.728 & 0.754 & 0.736 & 4,483 \\
Weighted Avg & 0.761 & 0.746 & 0.749 & 4,483 \\
\bottomrule
\end{tabular}
\end{table}

\textbf{Key Observations:}
\begin{itemize}
\item \textbf{Best performing:} Drop (90.7\% precision, 74.6\% recall) --- visually distinctive net-clearing motion
\item \textbf{Clear:} High recall (89.2\%) but lower precision (67.6\%) --- often confused with Lift
\item \textbf{Most challenging:} Drive (58.3\% precision, 61.6\% recall) --- fast flat motion similar to Lift/Smash
\end{itemize}

\subsubsection{Confusion Matrix Analysis}

\begin{table}[H]
\centering
\caption{Confusion Matrix (Test Set)}
\begin{tabular}{l|ccccc}
\toprule
True $\backslash$ Pred & Clear & Drive & Drop & Lift & Smash \\
\midrule
Clear & \textbf{477} & 26 & 5 & 20 & 7 \\
Drive & 37 & \textbf{456} & 47 & 145 & 55 \\
Drop & 117 & 85 & \textbf{1132} & 114 & 70 \\
Lift & 39 & 126 & 33 & \textbf{731} & 37 \\
Smash & 36 & 89 & 31 & 19 & \textbf{549} \\
\bottomrule
\end{tabular}
\end{table}

\textbf{Common Confusions:}
\begin{enumerate}
\item Drive $\leftrightarrow$ Lift (145 + 126 = 271 errors): Both are horizontal/defensive shots
\item Drop $\leftrightarrow$ Clear (117 + 5 = 122 errors): Both go to back of court
\item Smash $\leftrightarrow$ Drive (89 + 55 = 144 errors): Both are attacking shots with speed
\end{enumerate}

\subsection{Training Dynamics}

\begin{table}[H]
\centering
\caption{Training Progress}
\begin{tabular}{lcccc}
\toprule
Epoch & Train Acc & Val Acc & Train Loss & Val Loss \\
\midrule
10 & 61.23\% & 58.47\% & 1.124 & 1.213 \\
20 & 72.45\% & 69.32\% & 0.842 & 0.947 \\
30 & 79.12\% & 73.18\% & 0.651 & 0.812 \\
40 & 82.91\% & 74.91\% & 0.531 & 0.776 \\
44 & 84.03\% & 74.05\% & 0.498 & 0.789 \\
\bottomrule
\end{tabular}
\end{table}

\textbf{Observations:}
\begin{itemize}
\item Steady improvement through epoch 40
\item Train-val gap: 84.03\% - 74.05\% = 9.98\% (acceptable, no severe overfitting)
\item Early stopping at epoch 44 (validation accuracy plateaued)
\item Training time: ~6 hours on Mac M1 MPS
\end{itemize}

\section{Phase 3: Shot Coach Application}

\subsection{System Design}

Built a practical web application deploying the trained model for end-user shot analysis.

\subsubsection{Architecture}

\begin{figure}[H]
\centering
\begin{verbatim}
User uploads video (MP4/AVI/MOV)
    ↓
Shot Classifier Module
    → ResNet18+BiLSTM inference
    → Extract 16 frames uniformly
    → Return shot type + confidence + metadata
    ↓
Video Metadata Analyzer
    → Calculate video duration, FPS
    → Track sampled frame indices/times
    → Generate warnings for multi-shot videos
    ↓
Streamlit Web Interface
    → Display classification results
    → Show confidence scores (all 5 classes)
    → Provide video analysis metadata
    → Issue smart warnings
\end{verbatim}
\caption{Shot Coach Application Pipeline}
\end{figure}

\subsubsection{Key Features}

\textbf{1. Shot Classification:}
\begin{itemize}
\item Upload video via web interface
\item Automatic shot type detection (5 classes)
\item Confidence scores for all shot types
\item Processing time: 5-10 seconds per video
\end{itemize}

\textbf{2. Video Metadata Display:}
\begin{itemize}
\item Video duration, FPS, total frames
\item Frames analyzed by model (16 uniformly sampled)
\item Exact timeframe analyzed (start--end timestamps)
\item Frame sampling timeline visualization
\end{itemize}

\textbf{3. Smart Warning System:}

Implemented context-aware warnings based on video characteristics and confidence scores:

\begin{table}[H]
\centering
\caption{Warning Triggers}
\begin{tabular}{llp{6cm}}
\toprule
Condition & Warning Level & Message \\
\midrule
Duration $>$ 5s & ⚠️ Warning & Multi-shot video detected, may cause confused predictions \\
Duration $<$ 1.5s & ℹ️ Info & Too short, may not capture full motion \\
Confidence $<$ 60\% & ⚠️ Warning & Low confidence, check camera angle and lighting \\
Confidence $<$ 70\% \& Duration $>$ 3s & ⚠️ Warning & Low confidence with long video, likely multiple shots \\
Confidence $>$ 85\% & ✅ Success & High confidence, good video quality \\
\bottomrule
\end{tabular}
\end{table}

\subsection{Video Metadata Feature}

\subsubsection{Motivation}

Users asked two critical questions:
\begin{enumerate}
\item "If my video has many shots, which shot is analyzed?"
\item "Can you show what timeframe was analyzed?"
\end{enumerate}

These questions revealed a fundamental UX issue: users didn't understand that the model samples frames uniformly across the \textit{entire video}.

\subsubsection{Implementation}

Added comprehensive metadata tracking to the shot classifier:

\begin{lstlisting}[language=Python, caption=Metadata Return Format]
{
    'success': True,
    'predicted_class': 'Smash',
    'confidence': 0.943,
    'probabilities': {...},
    'metadata': {
        'total_frames': 90,
        'fps': 30.0,
        'duration': 3.0,
        'sampled_frames': [0, 5, 11, ..., 89],
        'sampled_times': [0.0, 0.17, 0.37, ..., 2.97],
        'num_frames_analyzed': 16
    }
}
\end{lstlisting}

\textbf{Benefits:}
\begin{enumerate}
\item \textbf{Transparency:} Users see exactly what was analyzed
\item \textbf{Education:} Explains uniform sampling strategy
\item \textbf{Better videos:} Users learn to trim to single shots
\item \textbf{Debugging:} Low confidence + long duration = likely multi-shot video
\end{enumerate}

\subsection{Camera Angle Robustness Analysis}

Investigated model performance sensitivity to camera viewpoint changes.

\subsubsection{Training Data Characteristics}

ShuttleSet videos recorded from professional broadcast cameras:
\begin{itemize}
\item \textbf{Angle:} Side-court view, elevated 45$^\circ$ from court level
\item \textbf{Distance:} 15--20 meters from court center
\item \textbf{Height:} 2--4 meters above court
\item \textbf{Consistency:} All 22,302 videos use similar camera positioning
\end{itemize}

\subsubsection{Expected Angle Sensitivity}

\begin{table}[H]
\centering
\caption{Predicted Performance by Camera Angle}
\begin{tabular}{lcc}
\toprule
Camera Angle & Expected Accuracy & Reason \\
\midrule
Side-court (training angle) & ✅ 74.6\% & Matches training distribution \\
Slight elevation change & ✅ 70--75\% & ResNet features rotation-invariant \\
Front/back court view & ⚠️ 50--65\% & Different perspective, motion patterns \\
Player POV (GoPro) & ❌ 30--50\% & Radically different, close-up \\
Overhead view & ❌ 30--50\% & Completely different, no depth \\
Close-up (upper body only) & ❌ 40--55\% & Missing footwork cues \\
\bottomrule
\end{tabular}
\end{table}

\textbf{Shot-Specific Vulnerabilities:}

\begin{itemize}
\item \textbf{Smash:} Low sensitivity (distinct downward motion visible from most angles)
\item \textbf{Clear:} Low sensitivity (high trajectory recognizable)
\item \textbf{Drop:} High sensitivity (subtle motion, hard to see from wrong angle)
\item \textbf{Drive:} Very high sensitivity (horizontal motion, easily confused from front view)
\item \textbf{Lift:} Moderate sensitivity (trajectory visible but may confuse with Clear)
\end{itemize}

\subsubsection{User Guidance}

Application provides camera angle recommendations:

\begin{itemize}
\item Record from side-court view (broadcast angle)
\item Position camera 3--5 meters away
\item Elevate camera 1--2 meters above court
\item Ensure full body visible in frame
\item Maintain consistent distance
\item Avoid: player POV, overhead, close-ups, extreme angles
\end{itemize}

\subsection{Limitations and Future Work}

\textbf{Current Limitations:}
\begin{enumerate}
\item Single-shot assumption (multi-shot videos cause confusion)
\item 2D analysis only (no depth/3D information)
\item Camera angle sensitivity (trained on broadcast views)
\item No automatic shot segmentation
\item No visual overlays or annotations
\end{enumerate}

\textbf{Future Enhancements:}
\begin{enumerate}
\item \textbf{Automatic shot detection:} Identify multiple shots in long videos using motion analysis
\item \textbf{Shot segmentation:} Split videos into individual shot clips automatically
\item \textbf{Multi-angle training:} Augment dataset with synthetic viewpoint changes
\item \textbf{Visual feedback:} Overlay skeleton, trajectory, and analysis on video
\item \textbf{Technique analysis:} Biomechanics scoring for form improvement (currently disabled due to pose detection issues)
\end{enumerate}

\section{Comparative Analysis}

\subsection{Pose-Only vs Video-Based}

\begin{table}[H]
\centering
\caption{Approach Comparison}
\begin{tabular}{lcc}
\toprule
Aspect & Pose-Only & Video-Based \\
\midrule
\textbf{Accuracy} & 50.0\% (random) & 74.6\% \\
\textbf{Classes} & 2 (Clear, Smash) & 5 (all types) \\
\textbf{Features} & 60 engineered & 512 learned (CNN) \\
\textbf{Input} & Skeletal keypoints & Raw pixels \\
\textbf{Model Params} & ~500K & ~12.6M \\
\textbf{Inference Time} & 5 ms & 150 ms \\
\textbf{Data Size} & 4,655 clips & 22,302 clips \\
\textbf{Preprocessing} & MediaPipe pose & Frame sampling \\
\textbf{Success} & ❌ Failed & ✅ Succeeded \\
\bottomrule
\end{tabular}
\end{table}

\subsection{Key Takeaways}

\begin{enumerate}
\item \textbf{Visual appearance dominates:} Racket orientation, ball trajectory, court context encode more information than body pose alone

\item \textbf{Equipment tracking essential:} Sports analysis requires tracking sports equipment, not just human body

\item \textbf{Transfer learning effective:} Pretrained ImageNet features generalize well to badminton domain

\item \textbf{End-to-end learning superior:} Learned features (512-dim CNN) outperform hand-crafted features (60-dim engineered)

\item \textbf{Temporal modeling matters:} BiLSTM captures motion patterns across time, critical for action recognition

\item \textbf{Class imbalance manageable:} Focal Loss effectively handles severe imbalance (3.5\% to 31.8\% class distribution)
\end{enumerate}

\section{AI Governance and Ethics}

Evaluated application against Singapore's AI Verify Framework \cite{AIVerify}.

\subsection{Transparency and Explainability}

\textbf{Implemented:}
\begin{itemize}
\item Clear display of model confidence scores (0--100\%)
\item Probabilities for all 5 shot types shown
\item Video metadata (what frames were analyzed)
\item Model architecture disclosed (ResNet18+BiLSTM)
\item Training accuracy reported (74.6\%)
\end{itemize}

\textbf{Limitations disclosed:}
\begin{itemize}
\item Camera angle sensitivity explained
\item Single-shot assumption stated
\item Multi-shot video warnings provided
\item Accuracy not 100\%, users warned
\end{itemize}

\subsection{Privacy and Data Protection}

\textbf{Design Principles:}
\begin{itemize}
\item \textbf{On-device processing:} Videos processed locally, not uploaded to servers
\item \textbf{No data retention:} Uploaded videos deleted after analysis
\item \textbf{No user tracking:} No analytics, cookies, or personal data collection
\item \textbf{No cloud dependencies:} Self-contained Streamlit app
\end{itemize}

\subsection{Safety and Robustness}

\textbf{Implemented Safeguards:}
\begin{enumerate}
\item \textbf{Confidence thresholds:} Warn users when confidence $<$ 70\%
\item \textbf{Input validation:} Check video format, duration, file size
\item \textbf{Graceful degradation:} Handle corrupted videos without crashing
\item \textbf{Error messaging:} Clear, actionable error messages
\end{enumerate}

\textbf{Disclaimers Required for Public Release:}
\begin{itemize}
\item "AI analysis for educational purposes only, not professional coaching advice"
\item "Accuracy 74.6\%, errors possible, verify critical decisions with coach"
\item "Optimized for side-court camera angles, other angles may reduce accuracy"
\item "Single shot per video required for accurate classification"
\end{itemize}

\subsection{Fairness and Inclusivity}

\textbf{Potential Biases:}
\begin{itemize}
\item Training data: Professional players only (high skill level)
\item Camera angles: Broadcast views (may not match recreational settings)
\item Shot types: International standard taxonomy (may differ by region)
\end{itemize}

\textbf{Mitigation Strategies:}
\begin{itemize}
\item Document training data characteristics
\item Provide camera angle guidance
\item Allow user feedback on incorrect classifications
\item Plan for diverse dataset expansion (recreational players, various angles)
\end{itemize}

\section{Conclusion}

This project demonstrated a complete machine learning development cycle from initial failure to production deployment.

\subsection{Key Achievements}

\begin{enumerate}
\item \textbf{Systematic failure analysis:} Rigorously investigated pose-only approach, identified fundamental limitations through statistical analysis

\item \textbf{Successful pivot:} Transitioned to video-based approach, achieved 74.6\% accuracy on 5-class classification (24.6 pp improvement over pose-only)

\item \textbf{Production application:} Deployed Shot Coach web app with video metadata tracking, smart warnings, and user-friendly interface

\item \textbf{Research contributions:}
   \begin{itemize}
   \item Empirical evidence that body pose alone insufficient for overhead shot discrimination
   \item Demonstration of transfer learning effectiveness for badminton action recognition
   \item Analysis of camera angle sensitivity for video-based shot classification
   \end{itemize}
\end{enumerate}

\subsection{Lessons Learned}

\begin{enumerate}
\item \textbf{Domain expertise early:} Understanding sport biomechanics before modeling saves months of effort

\item \textbf{Feature analysis first:} Analyze feature distributions before training complex models

\item \textbf{Fail fast:} Early negative results (50\% accuracy) enabled quick pivot to better approach

\item \textbf{User feedback critical:} Questions like "which shot is analyzed?" revealed UX gaps

\item \textbf{Transparency builds trust:} Showing metadata and warnings improves user confidence

\item \textbf{Practical deployment matters:} Technical accuracy (74.6\%) less important than usable application
\end{enumerate}

\subsection{Project Impact}

\textbf{Technical Impact:}
\begin{itemize}
\item Demonstrated video-based superiority over pose-only for badminton shot recognition
\item Established baseline (74.6\%) for future work
\item Open-sourced implementation for research community
\end{itemize}

\textbf{Practical Impact:}
\begin{itemize}
\item Functional tool for recreational players
\item Instant feedback without coach scheduling
\item Educational value in understanding shot types
\item Foundation for commercial coaching app
\end{itemize}

\subsection{Future Directions}

\textbf{Short-Term (3--6 months):}
\begin{enumerate}
\item Automatic shot segmentation in multi-shot videos
\item Visual overlays (skeleton, trajectory, analysis markers)
\item Mobile app version (iOS/Android)
\item Batch processing for practice session analysis
\end{enumerate}

\textbf{Medium-Term (6--12 months):}
\begin{enumerate}
\item Multi-angle training data augmentation
\item Shot quality scoring (success prediction)
\item Rally analysis (sequence of shots)
\item Comparison with professional players
\end{enumerate}

\textbf{Long-Term (1+ years):}
\begin{enumerate}
\item Real-time analysis (camera feed)
\item 3D pose estimation for depth information
\item Technique analysis with biomechanics feedback (revival of Phase 1 concepts with better data)
\item Personalized training recommendations based on shot patterns
\end{enumerate}

\subsection{Final Remarks}

This project exemplifies the iterative nature of applied machine learning: initial approaches fail, require re-thinking, pivot to alternative methods, and eventually produce practical solutions. The journey from 50\% accuracy (pose-only) to 74.6\% accuracy (video-based) to a deployed web application demonstrates that success requires not just technical skills, but domain understanding, user empathy, and willingness to abandon failed approaches.

The Shot Coach application, while limited, provides genuine value to recreational badminton players. It makes AI-powered shot analysis accessible through a simple web interface, educates users about video quality requirements through smart warnings, and establishes a foundation for future enhancements.

Most importantly, this project demonstrates that \textbf{AI governance starts at design time}---transparency, privacy, safety, and fairness were incorporated from the beginning, not retrofitted after deployment.

\section*{Acknowledgments}

This project used the ShuttleSet dataset \cite{ShuttleSet}, MediaPipe Pose from Google \cite{MediaPipe}, and PyTorch deep learning framework \cite{PyTorch}. Code development, debugging, report drafting, and LaTeX formatting were assisted by Claude Code \cite{ClaudeCode} and ChatGPT \cite{ChatGPT}. Special thanks to ITI123 instructors for guidance on AI governance principles.

\begin{thebibliography}{9}

\bibitem{ShuttleSet}
Wei-Yao Wang, Hong-Han Shuai, and Kai-Lung Hua.
\textit{ShuttleSet: A Human-Annotated Stroke-Level Singles Dataset for Badminton Tactical Analysis}.
CoRR, abs/2306.04948, 2023.
\url{https://arxiv.org/abs/2306.04948}

\bibitem{MediaPipe}
Google Research.
\textit{MediaPipe Pose: ML Solutions for Human Pose Estimation}.
\url{https://google.github.io/mediapipe/solutions/pose.html}, 2023.

\bibitem{FocalLoss}
Tsung-Yi Lin, Priya Goyal, Ross Girshick, Kaiming He, and Piotr Dollár.
\textit{Focal Loss for Dense Object Detection}.
IEEE International Conference on Computer Vision (ICCV), 2017.

\bibitem{PyTorch}
Adam Paszke et al.
\textit{PyTorch: An Imperative Style, High-Performance Deep Learning Library}.
Neural Information Processing Systems (NeurIPS), 2019.

\bibitem{AIVerify}
Infocomm Media Development Authority (IMDA), Singapore.
\textit{AI Verify: AI Governance Testing Framework and Toolkit}.
\url{https://aiverifyfoundation.sg}, 2023.

\bibitem{ClaudeCode}
Anthropic.
\textit{Claude Code: AI-Powered Command Line Coding Tool}.
\url{https://www.anthropic.com/claude/code}, 2025.

\bibitem{ChatGPT}
OpenAI.
\textit{ChatGPT: Optimizing Language Models for Dialogue}.
\url{https://openai.com/chatgpt}, 2024.

\end{thebibliography}

\newpage
\appendix

\section{Code Repository Structure}

\begin{lstlisting}[language=bash]
iti123_v2/
├── data/
│   ├── clips/          # 22,302 video clips (5 shot types)
│   ├── poses/          # MediaPipe pose extractions (Phase 1)
│   └── metadata.csv    # Shot annotations
├── notebooks/
│   ├── badminton_action_recognition_training.ipynb  # Phase 1 pose-only
│   └── badminton_video_classification.py            # Phase 2 video-based
├── scripts/
│   ├── extract_shuttleset_clips.py    # Video clip extraction
│   └── extract_poses_parallel.py      # Pose extraction
├── shot_coach/
│   ├── app.py                          # Streamlit web app
│   ├── modules/
│   │   ├── shot_classifier.py         # ResNet18+BiLSTM inference
│   │   ├── pose_extractor.py          # MediaPipe wrapper
│   │   └── technique_analyzer.py      # Biomechanics (disabled)
│   ├── test_shot_coach.py             # CLI testing
│   ├── requirements.txt               # Dependencies
│   └── README_CLASSIFIER.md           # Documentation
├── outputs/
│   ├── results_optionA/
│   │   ├── best_model.pth             # Trained model (164 MB)
│   │   ├── results_summary.json       # Training metrics
│   │   └── classification_report.txt  # Per-class performance
│   └── reports/
│       ├── ITI123_Milestone_Report.tex      # Phase 1 report
│       └── ITI123_Final_Project_Report.tex  # This document
└── docs/
    ├── VIDEO_TRAINING_README.md       # Training guide
    └── COLAB_SETUP_GUIDE.md          # Cloud training setup
\end{lstlisting}

\section{Model Training Commands}

\begin{lstlisting}[language=bash, caption=Phase 1: Pose-Only Training]
# Extract poses from videos
python scripts/extract_poses_parallel.py \
    --videos data/clips \
    --output data/poses \
    --workers 4

# Train pose-based models
cd notebooks
jupyter notebook badminton_action_recognition_training.ipynb
# Run all cells
# Result: 50% accuracy (failed)
\end{lstlisting}

\begin{lstlisting}[language=bash, caption=Phase 2: Video-Based Training]
# Train ResNet18+BiLSTM on videos
cd notebooks
python badminton_video_classification.py

# Expected output:
# Epoch 44/50
# Train Loss: 0.498 | Train Acc: 84.03%
# Val Loss: 0.789 | Val Acc: 74.05%
# Test Accuracy: 74.62%
# Model saved: outputs/results_optionA/best_model.pth
\end{lstlisting}

\begin{lstlisting}[language=bash, caption=Phase 3: Shot Coach Deployment]
# Install dependencies
cd shot_coach
pip install -r requirements.txt

# Run web application
streamlit run app.py

# Test with CLI
python test_shot_coach.py ../data/clips/Smash/video.mp4
\end{lstlisting}

\section{Sample Output}

\subsection{Shot Coach Classification Output}

\begin{verbatim}
======================================================================
SHOT COACH - CLASSIFICATION RESULTS
======================================================================

Detected Shot: Smash
Confidence: 94.3%

Video Analysis Info:
  Duration: 2.47s
  Total frames: 74
  Frames analyzed: 16
  Timeframe: 0.00s - 2.43s

All Shot Type Probabilities:
  🎯 Smash (Predicted)     ████████████████████████ 94.3%
  Clear                    ████░░░░░░░░░░░░░░░░░░░  4.2%
  Drive                    ░░░░░░░░░░░░░░░░░░░░░░░  1.1%
  Drop                     ░░░░░░░░░░░░░░░░░░░░░░░  0.3%
  Lift                     ░░░░░░░░░░░░░░░░░░░░░░░  0.1%

✅ High confidence - video angle matches training data well!

Shot Type Information:
💥 Smash - Powerful attacking shot hit downward. The primary
attacking shot in badminton.
======================================================================
\end{verbatim}

\end{document}
